% GENERATED FILE. Edit canonical lessons, then run npm run generate:books.
% canonical-source-hash: fnv1a64:18598df4904e96f9
% canonical-lessons: KA-C177-nikharavagi, KA-C177-sahajavagi, KA-C177-phalitamsa, KA-C177-parinama, KA-C177-vyatyasa

\chapter{Exactly, Naturally, Result, Effect, Difference}
\label{ch:ka-exactly-naturally-result-effect-difference}
\hlchaptermodality{177}

\begin{chapteropening}
\textbf{By the end of this chapter:} \emph{I can say exactly, naturally, a result, an effect and a difference.}

Complete the last lesson of chapter 177: I can say exactly, naturally, a result, an effect and a difference.
\end{chapteropening}

\section[\texorpdfstring{\kn{ನಿಖರವಾಗಿ} (nikharavāgi)}{ನಿಖರವಾಗಿ (nikharavāgi)}]{\kn{ನಿಖರವಾಗಿ} (nikharavāgi) — exactly, precisely}

\label{lesson:KA-C177-nikharavagi}



\begin{quote}
Before the new one: say the Kannada for in case, then the Kannada for instead.
\end{quote}

\subsection*{You'll want to know: \kn{ನಿಖರವಾಗಿ}}
\textbf{\kn{ನಿಖರವಾಗಿ}} — \emph{nikharavāgi} — \textquotedblleft{}exactly, precisely\textquotedblright{}.

\begin{grammarlens}[title={What you've built}]
One more word for this chapter.
\end{grammarlens}

\subsection*{Your turn}
\begin{itemize}
\raggedright
  \item \emph{Say it:} \emph{nikharavāgi}
  \item \emph{Say it:} \emph{nikharavāgi}, once more
  \item \emph{Say it:} say \emph{nikharavāgi}
  \item \emph{Recall:} say the Kannada for in case, then the Kannada for instead, then say \emph{nikharavāgi} again
\end{itemize}

\begin{tcolorbox}[breakable,colback=teal!4,colframe=teal!35!black,title={Before you move on}]
What does \kn{ನಿಖರವಾಗಿ} mean? (\textquotedblleft{}Exactly, precisely\textquotedblright{}.) Say it once more.
\end{tcolorbox}
\section[\texorpdfstring{\kn{ಸಹಜವಾಗಿ} (sahajavāgi)}{ಸಹಜವಾಗಿ (sahajavāgi)}]{\kn{ಸಹಜವಾಗಿ} (sahajavāgi) — naturally, of course}

\label{lesson:KA-C177-sahajavagi}



\begin{quote}
Before the new one: say the Kannada for instead, then the Kannada for exactly.
\end{quote}

\subsection*{You'll want to know: \kn{ಸಹಜವಾಗಿ}}
\textbf{\kn{ಸಹಜವಾಗಿ}} — \emph{sahajavāgi} — \textquotedblleft{}naturally, of course\textquotedblright{}.

\begin{grammarlens}[title={What you've built}]
One more word for this chapter.
\end{grammarlens}

\subsection*{Your turn}
\begin{itemize}
\raggedright
  \item \emph{Say it:} \emph{sahajavāgi}
  \item \emph{Say it:} \emph{sahajavāgi}, once more
  \item \emph{Say it:} say \emph{sahajavāgi}
  \item \emph{Recall:} say the Kannada for instead, then the Kannada for exactly, then say \emph{sahajavāgi} again
\end{itemize}

\begin{tcolorbox}[breakable,colback=teal!4,colframe=teal!35!black,title={Before you move on}]
What does \kn{ಸಹಜವಾಗಿ} mean? (\textquotedblleft{}Naturally, of course\textquotedblright{}.) Say it once more.
\end{tcolorbox}
\section[\texorpdfstring{\kn{ಫಲಿತಾಂಶ} (phalitāṃśa)}{ಫಲಿತಾಂಶ (phalitāṃśa)}]{\kn{ಫಲಿತಾಂಶ} (phalitāṃśa) — a result}

\label{lesson:KA-C177-phalitamsa}



\begin{quote}
Before the new one: say the Kannada for exactly, then the Kannada for naturally.
\end{quote}

\subsection*{You'll want to know: \kn{ಫಲಿತಾಂಶ}}
\textbf{\kn{ಫಲಿತಾಂಶ}} — \emph{phalitāṃśa} — \textquotedblleft{}a result\textquotedblright{}.

\begin{grammarlens}[title={What you've built}]
One more word for this chapter.
\end{grammarlens}

\subsection*{Your turn}
\begin{itemize}
\raggedright
  \item \emph{Say it:} \emph{phalitāṃśa}
  \item \emph{Say it:} \emph{phalitāṃśa}, once more
  \item \emph{Say it:} say \emph{phalitāṃśa}
  \item \emph{Recall:} say the Kannada for exactly, then the Kannada for naturally, then say \emph{phalitāṃśa} again
\end{itemize}

\begin{tcolorbox}[breakable,colback=teal!4,colframe=teal!35!black,title={Before you move on}]
What does \kn{ಫಲಿತಾಂಶ} mean? (\textquotedblleft{}A result\textquotedblright{}.) Say it once more.
\end{tcolorbox}
\section[\texorpdfstring{\kn{ಪರಿಣಾಮ} (pariṇāma)}{ಪರಿಣಾಮ (pariṇāma)}]{\kn{ಪರಿಣಾಮ} (pariṇāma) — an effect, a consequence}

\label{lesson:KA-C177-parinama}



\begin{quote}
Before the new one: say the Kannada for naturally, then the Kannada for a result.
\end{quote}

\subsection*{You'll want to know: \kn{ಪರಿಣಾಮ}}
\textbf{\kn{ಪರಿಣಾಮ}} — \emph{pariṇāma} — \textquotedblleft{}an effect, a consequence\textquotedblright{}.

\begin{grammarlens}[title={What you've built}]
One more word for this chapter.
\end{grammarlens}

\subsection*{Your turn}
\begin{itemize}
\raggedright
  \item \emph{Say it:} \emph{pariṇāma}
  \item \emph{Say it:} \emph{pariṇāma}, once more
  \item \emph{Say it:} say \emph{pariṇāma}
  \item \emph{Recall:} say the Kannada for naturally, then the Kannada for a result, then say \emph{pariṇāma} again
\end{itemize}

\begin{tcolorbox}[breakable,colback=teal!4,colframe=teal!35!black,title={Before you move on}]
What does \kn{ಪರಿಣಾಮ} mean? (\textquotedblleft{}An effect, a consequence\textquotedblright{}.) Say it once more.
\end{tcolorbox}
\section[\texorpdfstring{\kn{ವ್ಯತ್ಯಾಸ} (vyatyāsa)}{ವ್ಯತ್ಯಾಸ (vyatyāsa)}]{\kn{ವ್ಯತ್ಯಾಸ} (vyatyāsa) — a difference}

\label{lesson:KA-C177-vyatyasa}



\begin{quote}
Before the new one: say the Kannada for a result, then the Kannada for an effect.
\end{quote}

\subsection*{You'll want to know: \kn{ವ್ಯತ್ಯಾಸ}}
\textbf{\kn{ವ್ಯತ್ಯಾಸ}} — \emph{vyatyāsa} — \textquotedblleft{}a difference\textquotedblright{}.

\begin{grammarlens}[title={What you've built}]
One more word for this chapter.
\end{grammarlens}

\subsection*{Your turn}
\begin{itemize}
\raggedright
  \item \emph{Say it:} \emph{vyatyāsa}
  \item \emph{Say it:} \emph{vyatyāsa}, once more
  \item \emph{Say it:} say \emph{vyatyāsa}
  \item \emph{Recall:} say the Kannada for a result, then the Kannada for an effect, then say \emph{vyatyāsa} again
\end{itemize}

\begin{tcolorbox}[breakable,colback=teal!4,colframe=teal!35!black,title={Before you move on}]
What does \kn{ವ್ಯತ್ಯಾಸ} mean? (\textquotedblleft{}A difference\textquotedblright{}.) Say it once more.
\end{tcolorbox}
